% TC26 — SOS ฝั่งนิสิตไม่ปิด dialog เมื่อล้มเหลว
\begin{testscenario}{26}{3}
  \tcmeta{TC26}{CS 2568/01}{ระบบเว็บแอปพลิเคชัน ``NUCARE''}{นายสมชาย ใจดี}{วันที่ 5 พฤศจิกายน พ.ศ.\ 2569}{1.0}
  \tcname{การแจ้งช่วยเหลือด่วน (SOS) ของนิสิตเมื่อส่งไม่สำเร็จ}
  \tcpre{มีบัญชีนิสิตทดสอบ · รัน \tccode{student-sos.test.ts}}
  \tcdesc{ตรวจว่าการส่ง SOS ที่ล้มเหลวไม่ถูกมองว่าสำเร็จ และมีทางโทร 1323 / เข้าถึงจากมือถือและแชท}

  \tccase{1}{helper ไม่กลืน error}{%
    \begin{tcsteps}
      \item รัน \tccode{npx tsx --test src/lib/student-sos.test.ts}
    \end{tcsteps}}{%
    สำเร็จคืน ok · ล้มเหลวคืนข้อความ}{%
    PASS 2}{Pass}

  \tccase{2}{dialog คงอยู่เมื่อส่งไม่สำเร็จ}{%
    \begin{tcsteps}
      \item เปิดยืนยัน SOS แล้วจำลอง API error
    \end{tcsteps}}{%
    ไม่ปิด dialog · แสดงข้อความและชี้ 1323}{%
    การ์ดใช้ \tccode{runStudentSosRequest} และปิดเฉพาะเมื่อ ok}{Pass}

  \tccase{3}{ทางเข้าช่วยเหลือบนมือถือและแชทวิกฤต}{%
    \begin{tcsteps}
      \item ตรวจ bottom nav มีลิงก์ \tccode{/student\#emergency-help}
      \item ตรวจห้องแชทเมื่อ crisisSuggested มีปุ่มแจ้งทีมและโทร 1323
    \end{tcsteps}}{%
    เข้าถึงได้โดยไม่ต้องออกจากบริบทวิกฤต}{%
    ปรับใน \tccode{student-mobile-bottom-nav} และ \tccode{chat/page}}{Pass}
\end{testscenario}
